\subsection{Purpose}
Internships are a vital bridge between academic learning and professional experience, providing students with opportunities to apply their skills and companies with access to emerging talent. However, finding the right match between students and internships can be a complex and inefficient process. The Students\&Companies (S\&C) platform addresses these challenges by providing an integrated system that connects university students with companies offering internships. S\&C provides a wide range of features offering a seamless experience for students from the search of an internship to its completion and for companies an easy way to advertise their internship projects. The platform also lets universities to monitor and support their students throughout their internship, overseeing its execution and supporting the involved parties, should problems arise.
\subsection{Scope}
The Students\&Companies platform provides its services to different users, identifiable into three main classes: students,  companies and universities. Each one of them has different roles and interacts with the system in various ways.
S\&C is aimed mainly at university students that want to take part in an internship program, as they can search for new internships suiting their preferences or let the system inform them when internships matching their profile are available. S\&C is also able to provide suggestions on how to improve their profile to make it more appealing to companies.
The platform also supports companies easing the internships’ management. In doing so S\&C can be used to create questionnaires to select the most suitable candidates for an internship. Similarly to the suggestions for students, the system can analyze the internship offer to identify eventual improvements that can be made.
Hopefully, every internship comes to an end without problems but if that is not the case, interns and companies can communicate through the platform possibly involving the intern’s alma mater as a supervisor which can also terminate the internship if the parties cannot resolve the controversy.

\subsection{Definitions, Acronyms, Abbreviations}
\subsubsection{Definitions}
\begin{itemize}
    \item \textbf{Recommendations}: Process in which the system informs students when an internship that might interest them becomes available and can inform companies about the availability of student CVs corresponding to their needs.
    \item \textbf{SheerID}\cite{sheerid}: Platform that offers an API to verify students by using their academic credentials.
    \item \textbf{Contract}: The contract represents the lifecycle of an internship application, from initiation to its conclusion.
\end{itemize}
\subsubsection{Acronyms}
\begin{itemize}
    \item \textbf{S\&C}: Students \& Companies
    \item \textbf{DAO}: Data Accessing Object
    \item \textbf{API}: Application Programming Interface
    \item \textbf{REST}: REpresentational State Transfer
\end{itemize}
\subsubsection{Abbreviations}
\begin{itemize}
    \item \textbf{R*}: functional requirements.
\end{itemize}
\subsection{Revision History}
\begin{itemize}
    \item 07/01/2025 - Version 1.0.
\end{itemize}
\subsection{Reference Documents}
\begin{itemize}
    \item Assignment document “Assignment RDD AY 2024-2025.
\end{itemize}
\subsection{Document Structure}
\textbf{Introduction}: Provides a brief description of the project, focusing on the reasons behind its development and the goals to be achieved, as previously explained in the RASD.
Architectural Design: This section describes the architecture at different levels of granularity. \\
\textbf{User Interface Design}: This part presents the wireframes that represent the platform’s user interface. \\
\textbf{Requirements Traceability}: Explains how the goals defined in the previous document are fulfilled through the interaction between the requirements and the components described in this document. \\
\textbf{Implementation, Integration, and Test Plan}: Provides details about the implementation, integration, and testing processes, useful for developers and platform producers. \\
\textbf{Effort Spent}: Shows the effort involved in writing this document. \\
\textbf{Bibliography}: Contains references to any documents and software used to write this paper.